\section{Failures, Aligned Visualizations, and Video}
\label{app:qualitative}

\subsection{Prediction and Execution at the Same Physical Time}

A qualitative comparison uses the same physical episode and query time for
all candidates. Rows identify methods or interventions; columns identify
capture time and the physical prediction/application times. Predicted and
actual states are aligned to those times, rather than to equally spaced
video frame indices. Camera, reference frame, units, object identity, and
visibility are stated in the caption. A pose overlay uses measured/predicted
geometry and the actual calibration, not an illustrative reconstruction
presented as a measurement.

The draft currently contains no accepted rollout cases. Source episodes,
prediction arrays, application timestamps, and figure exports are
\placeholder{MISSING: real aligned qualitative assets}. Empty statistical
or trajectory panels are not replaced by synthetic curves or imagined
successes. A conceptual diagram, when useful, is labeled as a schematic
and carries no measured performance implication.

\subsection{Case Selection and Failure Attribution}

Case selection is declared before a visual claim is made. The final set should
cover the main supported comparisons and failures that qualify their scope,
using a reproducible episode-selection rule or an explicitly illustrative
selection. The selection rule, eligible episode pool, and selected identifiers
are \Pending. A visually clear success cannot establish a general success
rate, and a chosen failure cannot establish its population frequency.

Event traces first identify what happened: approach missed a feasible contact,
contact occurred at the wrong time or orientation, a grasp was not maintained,
a release failed containment, a collision occurred, or adaptation followed an
external impact. Whether the cause was perception, prediction, feasibility,
action generation, or stale execution is tested against synchronized outputs
and the appropriate comparison. These categories are diagnostic questions,
not a measured taxonomy. Their frequencies and attribution evidence are
\placeholder{MISSING: annotated failure results}.

Representation-specific failures remain visible. Examples include a target
prediction that omits a relevant obstacle, a context representation that
lacks grasp orientation, or a prediction whose average error hides a contact
mistake. These are hypotheses for inspection, not claimed observations.
Actual counterexamples and negative results, if they alter the Q1--Q3
conclusion, are summarized in the main paper with the full aligned record here.

\subsection{Supplementary Media Index}

Submission media are provided as anonymous supplementary files. The index
records filenames, durations, speed relative to real time, associated
configuration, and the result/case they illustrate. Videos disclose cuts,
slow motion, resets, and human intervention where present. A caption links a
clip to a measured episode; the clip does not replace quantitative evidence
or an essential method description in the main paper.

\begin{table}[t]
  \centering
  \caption{Media/asset index to fill from actual exports. No videos are
  asserted as delivered by this draft.}
  \label{tab:app_media_index}
  \small
  \begin{tabularx}{\columnwidth}{@{}lXX@{}}
    \toprule
    Item & Required content & File/case reference \\
    \midrule
    Q1 diagnostic & Matched current state and different histories & \Pending \\
    Q2 comparison & Predicted/actual future and contact outcome & \Pending \\
    Q3 intervention & Fixed current input; declared future/path change & \Pending \\
    Negative case & Full context and relevant event trace & \Pending \\
    External/real & Native task, time scale, and interventions & \Pending \\
    \bottomrule
  \end{tabularx}
\end{table}

The final media audit checks author/institution identifiers, screen overlays,
file metadata, audible identity cues, caption consistency, and whether the
files open at the supplied path. The actual file inventory and completed audit
are \Pending. Supplementary references below identify the sources used in
this document; they are independent of the main paper's bibliography.
